\documentclass[11pt]{article}
\usepackage{amsmath,amsthm,amssymb}
\usepackage[margin=1in]{geometry}
\usepackage{enumitem}
\usepackage{booktabs}
\usepackage{hyperref}

\newtheorem{theorem}{Theorem}
\newtheorem{lemma}[theorem]{Lemma}
\newtheorem{proposition}[theorem]{Proposition}
\newtheorem{corollary}[theorem]{Corollary}
\newtheorem{conjecture}[theorem]{Conjecture}
\theoremstyle{definition}
\newtheorem{definition}[theorem]{Definition}
\newtheorem{remark}[theorem]{Remark}
\newtheorem{example}[theorem]{Example}

\title{Structural $j$-formula for $\lambda = (5, 2, 1^q)$, $q \ge 5$:\\
       the second $r = 2$ same-row family of size $|\hat\lambda| = 7$}
\author{Clio}
\date{13 May 2026 (late night, prove session)}

\begin{document}
\maketitle

\begin{abstract}
The afternoon paper \texttt{2026-05-13-4-3-1n-upper-bound.tex}
established the structural $j$-formula
$B^{(\lambda)}_{j_0} V_\lambda \subseteq E_1^-|_{V_\lambda}$
for $\lambda = (4, 3, 1^q)$ at $q \ge 5$.
We carry out the same template for the second $r = 2$ family of
the same total $|\hat\lambda| = 7$, namely $\lambda = (5, 2, 1^q)$.

\textbf{Theorem.} For every $q \ge 5$,
\[
   B^{(\lambda)}_{j_0} V_\lambda \;\subseteq\; E_1^-|_{V_\lambda}
   \quad\text{and}\quad
   \dim B^{(\lambda)}_{j_0} V_\lambda \;\le\; 4
   \;=\; f^{\lambda^{(\ge 2)}}
   \;=\; f^{(4, 1)}.
\]
In particular, $\Pi^{S_n}|_{V_\lambda} = 0$ unconditionally for every
$n \ge 12$ (i.e., $q \ge 5$).

\textbf{Proof strategy.} Four-branch Phase A decomposition
\[
   E_{n-1}^+|_{V_\lambda}
   = \Phi_A(V_{\mu_A}) \,\oplus\, \Phi_B(V_{\mu_B})
     \,\oplus\, \Phi_C(V_{\mu_C}) \,\oplus\, S_1,
\]
with $\mu_A = (4, 2, 1^{q-1})$, $\mu_B = (5, 1^q)$,
$\mu_C = (4, 1^{q+1})$, and $S_1$ the same-row part at row $1$ with
sub-shape $\nu_1 = (3, 2, 1^q)$. Each piece is handled separately:
\begin{itemize}[leftmargin=2em,topsep=0.2em,itemsep=0.1em]
\item $\mu_A = (4, 2, 1^{q-1})$: $E_1^-$-containment via
\textbf{SRD-421} at sub-shape size $|\mu_A| = q + 5 \ge 10$
(structural at $q \ge 5$).
\item $\mu_B = (5, 1^q)$: hook of first part $5$, $E_1^-$-containment
via the \textbf{Shift Lemma all-hooks} paper at $q \ge 5$.
\item $\mu_C = (4, 1^{q+1})$: hook of first part $4$,
$j_0(\mu_C) = q + 3 = j_0(\lambda)$. The four-branch reduction
hits $\mu_C$ at index $j_0(\lambda) - 1 = q + 2 = \ell(\mu_C)$, one
step \emph{before} the sharp index. The extra $R'^{(\mu_C)}_{\ell(\mu_C)}$
step starts with $(T_1{+}1)$, which kills the $E_1^-$-image given by
the Shift Lemma hook $j$-formula on $\mu_C$ at $q + 1 \ge 4$, i.e.\
$q \ge 3$. \emph{$\mu_C$-leakage}: the piece vanishes outright.
\item $S_1$: same-row part at row $1$ with sub-shape $\nu_1 =
(3, 2, 1^q)$. The general same-row-dies framework
(\textbf{SRD-421} Theorem~14) applies because $\ell(\nu_1) =
\ell(\lambda)$ and the $j$-formula on $\nu_1$ at sharp index holds
by the May-20/May-12-today paper at $|\nu_1| = q + 5 \ge 8$ (i.e.\
$q \ge 3$). Hence $P_{j_0, n-1} S_1 = 0$.
\end{itemize}
Upper bound: $\dim \le \dim B^{(\mu_A)} + \dim B^{(\mu_B)} + 0 + 0
= 3 + 1 + 0 + 0 = 4$. The $E_1^-$ containment propagates through every
piece because $T_1$ commutes with each outer factor $(T_j{+}1)$ for
$j \ge j_0 - 1 = q + 2 \ge 7 \ge 3$.

\textbf{Significance.} This is the second structural-$j$-formula
family of size $|\hat\lambda| = 7$ after $(4, 3, 1^q)$. Combined,
the two papers complete the $|\hat\lambda| = 7$ row of the
SRD ladder, both at the threshold predicted by the night
$T_1 = -1$ theorem ($q_0((5,2)) = q_0((4,3)) = 5$). Computational
verification gives $\dim B = 4$ at $q = 5$ and the expected
non-stable behaviour below.
\end{abstract}

\section{Setup}\label{sec:setup}

Conventions follow the afternoon paper
(\texttt{2026-05-13-4-3-1n-upper-bound.tex}, henceforth
\textbf{(4,3,1)-UB}) and the structural-skeleton papers cited
therein. $H_q(S_n)$ is the Iwahori--Hecke algebra over $\mathbb{Q}(q)$
with generators $T_1, \ldots, T_{n-1}$ obeying $(T_i - q)(T_i + 1) = 0$
and braid. $V_\lambda$ is the seminormal cell module for
$\lambda \vdash n$, with basis $\{v_T : T \in \mathrm{SYT}(\lambda)\}$
and asymmetric Murphy normalisation. $E_j^\pm$ are the $q$- and
$(-1)$-eigenspaces of $T_j$. For $1 \le k \le n$,
\[
   R'_k \;:=\; (T_{k-1}{+}1)(T_{k-2}{+}1)\cdots(T_1{+}1) \in H_q(S_n),
   \qquad R'_1 := 1.
\]
For $1 \le a \le b \le n$, $P_{a, b} := R'_a R'_{a+1} \cdots R'_b$
(empty product is $1$). The $j$-image:
\[
   B^{(\lambda)}_j V_\lambda \;:=\; P_{j, n} V_\lambda,
   \qquad j_0(\lambda) := \ell(\lambda) + 1.
\]

\paragraph{Throughout this paper.}
$\lambda := (5, 2, 1^q)$ with $q \ge 1$, $n := q + 7$,
$\ell := \ell(\lambda) = q + 2$, $j_0 := j_0(\lambda) = q + 3 = n - 4$.

\paragraph{Removable corners of $\lambda$.}
\[
   c_1 := (1, 5) \text{ (length-pres)},
   \quad
   c_2 := (2, 2) \text{ (length-pres)},
   \quad
   c_3 := c_*(\lambda) = (q+2, 1) \text{ (column-$1$ bottom)}.
\]
Length-preservation of $c_1$: $\lambda_2 = 2 < 5 = \lambda_1$. \checkmark
Length-preservation of $c_2$: $\lambda_3 = 1 < 2 = \lambda_2$. \checkmark

\paragraph{Same-row pairs of $\lambda$.}
A same-row pair at row $i$ requires $\lambda_i \ge 2$ and that the SYT
placement of $\{n-1, n\}$ at $(i, \lambda_i - 1), (i, \lambda_i)$ is
realisable without violating column-strict increase in column~$1$.

\emph{Row $1$.} $\lambda_1 = 5 \ge 2$, cells $(1, 4), (1, 5)$.
$\lambda_2 = 2$, so cell $(2, 4), (2, 5)$ don't exist; no conflict from
row $2$. Letters at cells in column~$\ge 4$ above letter $n-1$ are
$\le n - 2$; rest of the shape fills with $\{2, \ldots, n-2\}$. The
sub-shape $\nu_1 := \lambda \setminus \{(1,4), (1,5)\} = (3, 2, 1^q)$
is a valid partition. \emph{Same-row pair at row $1$ exists.}

\emph{Row $2$.} $\lambda_2 = 2$, cells $(2, 1), (2, 2)$. Letter at
$(2, 1) = n-1$ but then column~$1$ below requires entries $> n - 1$,
i.e.\ $= n$; but $n$ is at $(2, 2)$. So $(3, 1), \ldots, (q+2, 1)$ have
no valid letters when $q \ge 1$. \emph{No same-row pair at row $2$.}

Rows $\ge 3$ have $\lambda_i = 1$; no same-row pairs.

\paragraph{Inner partitions on $n - 2$ letters.}
For each corner pair $X = \{c, c'\}$ of $\lambda$, set $\mu_X
:= \lambda \setminus X$:
\begin{align*}
   \mu_A &:= \lambda \setminus \{c_1, c_3\} = (4, 2, 1^{q-1}),
      &(\text{top corner $c_1$ paired with bottom $c_3$})\\
   \mu_B &:= \lambda \setminus \{c_2, c_3\} = (5, 1^q),
      &(\text{middle corner $c_2$ paired with bottom $c_3$})\\
   \mu_C &:= \lambda \setminus \{c_1, c_2\} = (4, 1^{q+1}),
      &(\text{both length-pres corners paired}).
\end{align*}
All three have $|\mu_X| = q + 5 = n - 2$, and
\[
   \ell(\mu_A) = q + 1,
   \quad
   \ell(\mu_B) = q + 1,
   \quad
   \ell(\mu_C) = q + 2 = \ell(\lambda).
\]
Sharp indices: $j_0(\mu_A) = j_0(\mu_B) = q + 2 = j_0(\lambda) - 1$;
$j_0(\mu_C) = q + 3 = j_0(\lambda)$.

\paragraph{Same-row sub-shape.}
$\nu_1 := \lambda \setminus \{(1, 4), (1, 5)\} = (3, 2, 1^q)$, a partition
on $n - 2 = q + 5$ letters with $\ell(\nu_1) = q + 2 = \ell(\lambda)$
and $j_0(\nu_1) = q + 3 = j_0(\lambda)$.

\section{Known ingredients}\label{sec:ingredients}

We re-cite the toolkit, verbatim where possible from \textbf{(4,3,1)-UB}.

\begin{theorem}[Phase A endpoint, May-19]\label{thm:phaseA}
For any $H_q(S_n)$-module $V$, $R'_n V = E_{n-1}^+(V)$.
\end{theorem}

\begin{theorem}[Corner-pair embedding $\Phi_X$, May-20 + SRD-421]\label{thm:Phi}
For each corner pair $X = \{c, c'\}$ of $\lambda$ with $|d_X| \ge 2$
(axial distance), there is a $\mathbb{Q}(q)$-linear injection
$\Phi_X : V_{\mu_X} \hookrightarrow V_\lambda$ such that:
\begin{enumerate}[leftmargin=2em,topsep=0.2em,itemsep=0.1em]
\item For each SYT $T'$ of $\mu_X$, $\Phi_X(v_{T'})$ is the
$T_{n-1}$-$(+q)$-eigenvector of the $2$-block
$\{v_{T_A^X}, v_{T_B^X}\}$ where $T_A^X, T_B^X$ extend $T'$ by
placing $\{n-1, n\}$ at $X$ in opposite orders.
\item $T_i \, \Phi_X(v) = \Phi_X(T_i v)$ for every
$1 \le i \le n - 3$ and every $v \in V_{\mu_X}$.
\item $\Phi_X(V_{\mu_X})$ is the span of $T_{n-1}$-$(+q)$-eigenvectors
in the corner-pair $2$-blocks for $X$.
\item For $n \ge 4$ (so $1 \le n - 3$), $\Phi_X$ takes
$E_1^-|_{V_{\mu_X}}$ injectively into $E_1^-|_{V_\lambda}$.
\end{enumerate}
\end{theorem}

\paragraph{Verification of $|d_X| \ge 2$ for $\lambda = (5, 2, 1^q)$.}
\begin{itemize}[leftmargin=2em,topsep=0.2em,itemsep=0.1em]
\item $X = A = \{c_1, c_3\} = \{(1, 5), (q+2, 1)\}$:
$d_A = (1 - 5) - ((q+2) - 1) = -q - 5$, so $|d_A| = q + 5 \ge 6$.
\item $X = B = \{c_2, c_3\} = \{(2, 2), (q+2, 1)\}$:
$d_B = (2 - 2) - ((q+2) - 1) = -q - 1$, so $|d_B| = q + 1 \ge 2$.
\item $X = C = \{c_1, c_2\} = \{(1, 5), (2, 2)\}$:
$d_C = (1 - 5) - (2 - 2) = -4$, so $|d_C| = 4 \ge 2$.
\end{itemize}
All three pairs are non-degenerate; $\Phi_A, \Phi_B, \Phi_C$ exist.

\begin{lemma}[Phase A decomposition for $\lambda = (5, 2, 1^q)$]\label{lem:phaseA-decomp}
For $\lambda = (5, 2, 1^q)$ at every $q \ge 1$,
\begin{equation}\label{eq:phaseA-decomp}
   E_{n-1}^+|_{V_\lambda}
   \;=\;
   \Phi_A(V_{\mu_A})
   \,\oplus\, \Phi_B(V_{\mu_B})
   \,\oplus\, \Phi_C(V_{\mu_C})
   \,\oplus\, S_1,
\end{equation}
where $S_1 := \mathrm{span}\{v_T : T(1, 4) = n - 1,\, T(1, 5) = n\}$ is
the same-row part at row $1$.
\end{lemma}

\begin{proof}
$+q$-eigenvectors of $T_{n-1}$ in the seminormal basis of $V_\lambda$
arise from two configurations of $\{n-1, n\}$:
(i) at a corner pair $\{c, c'\}$ with $c \ne c'$ in different rows and
columns (non-degenerate $2$-block), giving the
$+q$-eigenvector $\Phi_{\{c, c'\}}(v_{T'})$ via Theorem~\ref{thm:Phi}(3);
or (ii) in a same-row pair, giving $v_T$ itself as the
$+q$-eigenvector (Hoefsmit: $T_{n-1} v_T = q v_T$ when $n-1, n$ are in
the same row).

For $\lambda = (5, 2, 1^q)$ the three corner pairs are $A, B, C$
above, and the unique same-row pair is at row $1$ (see Setup).
The four families have pairwise disjoint seminormal supports (each is
characterised by where $\{n-1, n\}$ sit), so the sum
\eqref{eq:phaseA-decomp} is direct. The total counts every
$+q$-eigenvector of $T_{n-1}$ in $V_\lambda$ exactly once.
\end{proof}

\begin{theorem}[$\mu_A$: SRD-421, structural at $|\mu_A| \ge 10$]\label{thm:muA}
For $\mu_A = (4, 2, 1^{q-1})$:
\begin{enumerate}[leftmargin=2em,topsep=0.2em,itemsep=0.1em]
\item $B^{(\mu_A)}_{j_0(\mu_A)} V_{\mu_A} \subseteq E_1^-|_{V_{\mu_A}}$
\emph{structurally} at $|\mu_A| \ge 10$ (i.e.\ $q - 1 \ge 4$,
equivalently $q \ge 5$).
\item $\dim B^{(\mu_A)}_{j_0(\mu_A)} V_{\mu_A} = 3$ \emph{structurally}
in the same range.
\end{enumerate}
\end{theorem}

\begin{theorem}[$\mu_B$: hook of first part $5$, structural at $q \ge 5$]\label{thm:muB}
For $\mu_B = (5, 1^q)$ (a hook with $|\hat\mu_B| = 5$,
$q(\mu_B) = q$):
\begin{enumerate}[leftmargin=2em,topsep=0.2em,itemsep=0.1em]
\item \emph{(Shift Lemma all-hooks.)} The $j$-formula
$B^{(\mu_B)}_{j_0(\mu_B)} V_{\mu_B} \subseteq E_1^-|_{V_{\mu_B}}$
holds for every $q(\mu_B) \ge 5$, i.e.\ $q \ge 5$ in our normalisation.
\item $\dim B^{(\mu_B)}_{j_0(\mu_B)} V_{\mu_B} = 1$ in the same range
(hook dim $= 1$).
\end{enumerate}
\end{theorem}

\begin{theorem}[$\mu_C$: hook of first part $4$, structural at $q \ge 3$]\label{thm:muC}
For $\mu_C = (4, 1^{q+1})$:
\[
   B^{(\mu_C)}_{j_0(\mu_C)} V_{\mu_C} \;\subseteq\; E_1^-|_{V_{\mu_C}}
   \quad\text{for every } q + 1 \ge 4,\; \text{i.e.\ } q \ge 3.
\]
\end{theorem}

\begin{theorem}[$\nu_1 = (3, 2, 1^q)$: May-20 + May-12-today]\label{thm:nu1}
For $\nu_1 = (3, 2, 1^q)$ at $|\nu_1| = q + 5 \ge 8$, i.e.\ $q \ge 3$,
\[
   B^{(\nu_1)}_{j_0(\nu_1)} V_{\nu_1} \;\subseteq\; E_1^-|_{V_{\nu_1}}.
\]
\end{theorem}

\begin{theorem}[General same-row-dies framework, SRD-421 Theorem~14]\label{thm:srd-framework}
Let $\widetilde\lambda \vdash N$ admit a same-row pair at row $i$, set
$\widetilde\nu_i := \widetilde\lambda \setminus \{(i, \widetilde\lambda_i - 1),
(i, \widetilde\lambda_i)\}$, $\widetilde\ell := \ell(\widetilde\lambda)$,
$\widetilde j_0 := \widetilde\ell + 1$, and $\widetilde S_i$ the
corresponding same-row span. Assume $\widetilde\lambda_i \ge 3$
(so $\ell(\widetilde\nu_i) = \widetilde\ell$) and
$B^{(\widetilde\nu_i)}_{j_0(\widetilde\nu_i)} V_{\widetilde\nu_i}
\subseteq E_1^-|_{V_{\widetilde\nu_i}}$. Then
\[
   R'_{\widetilde j_0}\, R'_{\widetilde j_0 + 1}\, \cdots\, R'_{N-1}
   \, \widetilde S_i \;=\; 0.
\]
\end{theorem}

\section{The four-branch reduction}\label{sec:reduction}

The reduction mirrors \textbf{(4,3,1)-UB}~\S3 with the same commutation
lemma and four-branch identity.

\begin{lemma}[Commutation lemma]\label{lem:commute}
For each $X \in \{A, B, C\}$ and every $j_0 \le k \le n - 1$,
$R'_k\, \Phi_X(v) = (T_{k-1}{+}1) \, \Phi_X(R'^{(\mu_X)}_{k-1} v)$ for
$v \in V_{\mu_X}$. Iterating,
\begin{equation}\label{eq:commute-P}
   P_{j_0, n-1}\, \Phi_X(V_{\mu_X})
   \;=\; (T_{j_0-1}{+}1)(T_{j_0}{+}1)\cdots(T_{n-2}{+}1)
       \, \Phi_X\!\big(P^{(\mu_X)}_{j_0-1, n-2}\, V_{\mu_X}\big).
\end{equation}
\end{lemma}

\begin{proof}
Verbatim from \textbf{(4,3,1)-UB}~Lemma~10. The proof uses only the
$\langle T_1, \ldots, T_{n-3}\rangle$-equivariance of $\Phi_X$
(Theorem~\ref{thm:Phi}(2)) and the index-gap commutation of $T_j$
with $T_i$ for $|j - i| \ge 2$ inside $H_q(S_n)$.
\end{proof}

\begin{theorem}[Four-branch reduction]\label{thm:reduction}
For $\lambda = (5, 2, 1^q)$ at every $q \ge 1$,
\begin{equation}\label{eq:reduction}
   B^{(\lambda)}_{j_0} V_\lambda
   \;=\;
   (T_{j_0-1}{+}1)(T_{j_0}{+}1)\cdots(T_{n-2}{+}1)
   \Bigl[
   \sum_{X \in \{A, B, C\}} \Phi_X\!\big(P^{(\mu_X)}_{j_0-1, n-2} V_{\mu_X}\big)
   \Bigr]
   \;+\; P_{j_0, n-1}\, S_1.
\end{equation}
\end{theorem}

\begin{proof}
$B^{(\lambda)}_{j_0} V_\lambda = P_{j_0, n-1}\, R'_n V_\lambda$ by
definition. By Theorem~\ref{thm:phaseA} and
Lemma~\ref{lem:phaseA-decomp},
$R'_n V_\lambda = E_{n-1}^+|_{V_\lambda} = \Phi_A(V_{\mu_A}) \oplus
\Phi_B(V_{\mu_B}) \oplus \Phi_C(V_{\mu_C}) \oplus S_1$. Apply
$P_{j_0, n-1}$ and use Lemma~\ref{lem:commute} on each $\Phi_X$-summand.
\end{proof}

\section{Each piece evaluated}\label{sec:pieces}

\subsection{$\mu_A = (4, 2, 1^{q-1})$ piece (structural at $q \ge 5$)}\label{ssec:muA-piece}

By Theorem~\ref{thm:muA}(1), at $q \ge 5$,
$B^{(\mu_A)}_{j_0(\mu_A)} V_{\mu_A} \subseteq E_1^-|_{V_{\mu_A}}$.
We have $j_0(\mu_A) = j_0 - 1$ and $|\mu_A| = n - 2$, so
\[
   P^{(\mu_A)}_{j_0 - 1, n - 2}\, V_{\mu_A}
   \;=\; B^{(\mu_A)}_{j_0(\mu_A)} V_{\mu_A}
   \;\subseteq\; E_1^-|_{V_{\mu_A}}.
\]
By Theorem~\ref{thm:Phi}(4),
$\Phi_A(P^{(\mu_A)}_{j_0 - 1, n - 2} V_{\mu_A}) \subseteq E_1^-|_{V_\lambda}$.

The outer factors $(T_{j_0-1}{+}1)\cdots(T_{n-2}{+}1)$ have indices
$\ge j_0 - 1 = q + 2 \ge 7$, in particular $\ge 3$, so they commute
with $T_1$ inside $H_q(S_n)$ and preserve $E_1^-|_{V_\lambda}$.
Therefore the $A$-piece of \eqref{eq:reduction} lies in
$E_1^-|_{V_\lambda}$. Its dimension is at most $\dim B^{(\mu_A)} = 3$.

\subsection{$\mu_B = (5, 1^q)$ piece (structural at $q \ge 5$)}\label{ssec:muB-piece}

By Theorem~\ref{thm:muB}(1), at $q \ge 5$,
$B^{(\mu_B)}_{j_0(\mu_B)} V_{\mu_B} \subseteq E_1^-|_{V_{\mu_B}}$.
The same reasoning as for the $A$-piece (using
$j_0(\mu_B) = j_0 - 1$ and the same outer factors) gives
$\Phi_B(P^{(\mu_B)}_{j_0 - 1, n - 2} V_{\mu_B}) \subseteq
E_1^-|_{V_\lambda}$, with dimension at most $\dim B^{(\mu_B)} = 1$.

\subsection{$\mu_C = (4, 1^{q+1})$ piece: leakage to zero
($q \ge 3$)}\label{ssec:muC-leakage}

Here $j_0(\mu_C) = q + 3 = j_0(\lambda)$, one step \emph{above}
$j_0 - 1$. So
\[
   P^{(\mu_C)}_{j_0 - 1, n - 2} V_{\mu_C}
   \;=\; R'^{(\mu_C)}_{j_0 - 1} \cdot P^{(\mu_C)}_{j_0, n-2} V_{\mu_C}
   \;=\; R'^{(\mu_C)}_{j_0 - 1} \cdot B^{(\mu_C)}_{j_0(\mu_C)} V_{\mu_C}.
\]
By Theorem~\ref{thm:muC} (hook $j$-formula at $q + 1 \ge 4$, i.e.,
$q \ge 3$), $B^{(\mu_C)}_{j_0(\mu_C)} V_{\mu_C} \subseteq
E_1^-|_{V_{\mu_C}}$. The operator $R'^{(\mu_C)}_{j_0 - 1} =
(T_{j_0-2}{+}1)(T_{j_0-3}{+}1)\cdots(T_1{+}1)$ ends with $(T_1{+}1)$,
which kills $E_1^-|_{V_{\mu_C}}$ (since $T_1 = -1$ there). Hence
\[
   P^{(\mu_C)}_{j_0 - 1, n - 2} V_{\mu_C} \;=\; 0,
\]
and the $C$-piece of \eqref{eq:reduction} vanishes outright.

\subsection{Same-row part $S_1$: SRD-framework + $\nu_1 = (3, 2, 1^q)$
($q \ge 3$)}\label{ssec:S1-dies}

The same-row pair lives at row~$1$ of $\lambda$ with sub-shape
$\nu_1 = (3, 2, 1^q)$. We check the hypotheses of
Theorem~\ref{thm:srd-framework} (\textbf{SRD-421} Theorem~14):
\begin{itemize}[leftmargin=2em,topsep=0.2em,itemsep=0.1em]
\item $\widetilde\lambda = \lambda$, $i = 1$, $\widetilde\lambda_i =
\lambda_1 = 5 \ge 3$. \checkmark
\item $\ell(\nu_1) = q + 2 = \ell(\lambda) = \widetilde\ell$,
$\widetilde j_0 = j_0(\lambda)$. \checkmark
\item $B^{(\nu_1)}_{j_0(\nu_1)} V_{\nu_1} \subseteq E_1^-|_{V_{\nu_1}}$
at $q \ge 3$ by Theorem~\ref{thm:nu1}. \checkmark
\end{itemize}
Therefore $R'_{j_0}\, R'_{j_0+1} \cdots R'_{n-1}\, S_1
= P_{j_0, n-1}\, S_1 = 0$.

\section{The main theorem}\label{sec:main}

\begin{theorem}[$j$-formula and upper bound on $\dim B$ for $(5, 2, 1^q)$]
\label{thm:main}
For $\lambda = (5, 2, 1^q)$ at every $q \ge 5$,
\[
   B^{(\lambda)}_{j_0} V_\lambda \;\subseteq\; E_1^-|_{V_\lambda},
   \qquad
   \dim B^{(\lambda)}_{j_0} V_\lambda \;\le\; 4 \;=\; f^{(4, 1)}.
\]
\end{theorem}

\begin{proof}
By Theorem~\ref{thm:reduction},
$B^{(\lambda)}_{j_0} V_\lambda = \mathcal{O}_\lambda
[\Phi_A(P^{(\mu_A)}_{j_0-1, n-2} V_{\mu_A}) + \Phi_B(P^{(\mu_B)}_{j_0-1, n-2}
V_{\mu_B}) + \Phi_C(P^{(\mu_C)}_{j_0-1, n-2} V_{\mu_C})] +
P_{j_0, n-1} S_1$, where $\mathcal{O}_\lambda = (T_{j_0-1}{+}1)\cdots
(T_{n-2}{+}1)$.

By Sections~\ref{ssec:muC-leakage} and~\ref{ssec:S1-dies}, the $C$-piece
and $S_1$-piece vanish at $q \ge 3$. By Sections~\ref{ssec:muA-piece}--%
\ref{ssec:muB-piece}, the $A$- and $B$-pieces (after the outer factor
$\mathcal{O}_\lambda$) lie in $E_1^-|_{V_\lambda}$ at $q \ge 5$, with
$\dim \le 3 + 1 = 4$.

So at $q \ge 5$, $B^{(\lambda)}_{j_0} V_\lambda \subseteq E_1^-|_{V_\lambda}$
and $\dim B^{(\lambda)}_{j_0} V_\lambda \le 4 = f^{(4,1)}$.
\end{proof}

\begin{corollary}[Rank-zero theorem for $(5, 2, 1^{n-7})$, $n \ge 12$]
\label{cor:rank-zero}
For every $n \ge 12$, $\Pi^{S_n}|_{V_\lambda} = 0$ where
$\lambda = (5, 2, 1^{n-7})$.
\end{corollary}

\begin{proof}
By Theorem~\ref{thm:main}, $B^{(\lambda)}_{j_0} V_\lambda \subseteq
E_1^-|_{V_\lambda}$. Then $\Pi^{S_n}|_{V_\lambda} = R'_2 \cdot R'_3 \cdots
R'_{j_0 - 1} \cdot B^{(\lambda)}_{j_0} V_\lambda$. The trailing factor
of $R'_2$ is $(T_1 + 1)$, which annihilates $E_1^-$. Hence
$\Pi^{S_n}|_{V_\lambda} = 0$.
\end{proof}

\section{Matching lower bound: $\dim B \ge 4$}\label{sec:LB}

The four-branch reduction supplies an upper bound $\dim B \le 4$. The
matching lower bound follows from the
\textbf{disjoint-cell-pair $\Phi$-supports + adjacent-injectivity}
template of the May-13 \textbf{LB-closed} paper, applied verbatim
here.

\begin{lemma}[Disjoint $\Phi_A, \Phi_B$ supports at the $V_\lambda$ level]\label{lem:disjoint-AB}
$\Phi_A(V_{\mu_A}) \cap \Phi_B(V_{\mu_B}) = 0$ in $V_\lambda$.
\end{lemma}

\begin{proof}
$\Phi_X(v_{T'})$ is supported on the two SYTs of $\lambda$ that
place $\{n-1, n\}$ at the cell pair $X$ in the two orders, with $T'$
filling the rest. The unordered cell pairs are
$X = A: \{c_1, c_3\} = \{(1, 5), (q+2, 1)\}$ and
$X = B: \{c_2, c_3\} = \{(2, 2), (q+2, 1)\}$. They share $c_3$ but
differ in the other element: $c_1 = (1, 5) \ne (2, 2) = c_2$.

Any SYT $T$ supporting $\Phi_A(v_{T'_A})$ has $\{n-1, n\}$ at
$\{(1, 5), (q+2, 1)\}$; any SYT $T$ supporting $\Phi_B(v_{T'_B})$
has $\{n-1, n\}$ at $\{(2, 2), (q+2, 1)\}$. These cell pairs are
distinct as unordered sets, so no SYT of $\lambda$ supports both
families. Hence the seminormal supports of $\Phi_A(V_{\mu_A})$ and
$\Phi_B(V_{\mu_B})$ are disjoint, and the intersection is $0$.
\end{proof}

\begin{theorem}[Adjacent-injectivity of $\mathcal{O}_\lambda$ on $E_{n-1}^+|_{V_\lambda}$,
May-13 dim2-331-structural]\label{thm:adjacent-inj}
For any $H_q(S_n)$-module $V$ and $a \le b \le n - 2$, the operator
$(T_a{+}1)(T_{a+1}{+}1)\cdots(T_b{+}1)$ is injective from
$E_{b+1}^+(V)$ into $E_a^+(V)$ provided each intermediate
$\pi_j^+ |_{E_{j+1}^+}$ is injective (May-19 adjacent-disjointness
lemma). In particular, the outer chain
$\mathcal{O}_\lambda = (T_{j_0-1}{+}1)\cdots(T_{n-2}{+}1)$ is injective
on $E_{n-1}^+|_{V_\lambda}$.
\end{theorem}

\begin{theorem}[$\dim B \ge 4$ structurally at $q \ge 5$]\label{thm:LB}
For $\lambda = (5, 2, 1^q)$ at every $q \ge 5$,
$\dim B^{(\lambda)}_{j_0} V_\lambda \ge 4$.
\end{theorem}

\begin{proof}
By the four-branch reduction (Theorem~\ref{thm:reduction}) and the
$C$-piece / $S_1$-piece vanishing (Sections~\ref{ssec:muC-leakage},
\ref{ssec:S1-dies}),
\[
   B^{(\lambda)} V_\lambda
   \;=\; \mathcal{O}_\lambda\!\Big[\,
      \Phi_A\!\big(B^{(\mu_A)}\big) \;+\; \Phi_B\!\big(B^{(\mu_B)}\big)
   \Big].
\]
Both $\Phi_A(B^{(\mu_A)})$ and $\Phi_B(B^{(\mu_B)})$ are subspaces of
$\Phi_A(V_{\mu_A})$ and $\Phi_B(V_{\mu_B})$ respectively; by
Lemma~\ref{lem:disjoint-AB}, these are linearly independent. Hence
\[
   \Phi_A(B^{(\mu_A)}) \;+\; \Phi_B(B^{(\mu_B)})
   \;=\; \Phi_A(B^{(\mu_A)}) \;\oplus\; \Phi_B(B^{(\mu_B)}),
\]
of total dimension $\dim B^{(\mu_A)} + \dim B^{(\mu_B)} = 3 + 1 = 4$
at $q \ge 5$.

Both summands lie in $E_{n-1}^+|_{V_\lambda}$ (Theorem~\ref{thm:Phi}(3)).
By Theorem~\ref{thm:adjacent-inj}, $\mathcal{O}_\lambda$ is injective on
$E_{n-1}^+|_{V_\lambda}$. Hence
\[
   \dim B^{(\lambda)} V_\lambda
   \;\ge\; \dim \big(\Phi_A(B^{(\mu_A)}) \oplus \Phi_B(B^{(\mu_B)})\big)
   \;=\; 4.
\]
\end{proof}

\begin{corollary}[$\dim B = 4 = f^{\lambda^{(\ge 2)}}$ structurally at $q \ge 5$]
\label{cor:dim4}
For $\lambda = (5, 2, 1^q)$ at every $q \ge 5$,
$\dim B^{(\lambda)}_{j_0} V_\lambda = 4 = f^{(4, 1)} =
f^{\lambda^{(\ge 2)}}$.
\end{corollary}

\begin{proof}
Combine Theorem~\ref{thm:main} (upper bound) and Theorem~\ref{thm:LB}
(lower bound).
\end{proof}

\begin{remark}[Match with night $T_1 = -1$ threshold]\label{rem:match-night}
The night paper \texttt{2026-05-13-night-T1-neg1-deep-stable.tex} gave
the threshold $q_0(\hat\lambda) = |\hat\lambda| - 2\ell(\hat\lambda) + 2$
for $B \subseteq E_1^-$ under the fully-stable hypothesis. For
$\hat\lambda = (5, 2)$: $q_0 = 7 - 4 + 2 = 5$. The present theorem
realises this prediction structurally and unconditionally (the
fully-stable hypothesis is discharged by the SRD-framework chain
inside the proof).

Compare with $(4, 3, 1^q)$: $\hat\lambda = (4, 3)$, $|\hat\lambda| =
7$, $\ell(\hat\lambda) = 2$, $q_0 = 5$. Same threshold. The two
$|\hat\lambda| = 7$ families coincide on threshold and on
$f^{\lambda^{(\ge 2)}}$: for $(4, 3, 1^q)$, $\lambda^{(\ge 2)} = (3, 2)$,
$f = 5$; for $(5, 2, 1^q)$, $\lambda^{(\ge 2)} = (4, 1)$, $f = 4$.
Differential dim, identical threshold.
\end{remark}

\section{Computational verification}\label{sec:verify}

Mod-$P$ arithmetic at $P = 100\,003$ and $q = 11$. For each test shape,
build the matrix $M = R'_{j_0} R'_{j_0+1} \cdots R'_n$ via Hoefsmit
seminormal matrices for $T_1, \ldots, T_{n-1}$, then check $\dim B =
\mathrm{rank}(M)$ and $(T_1 + 1) \cdot B = 0$.

\paragraph{$\lambda = (5, 2, 1^q)$.}
\begin{center}
\begin{tabular}{ccccc}
\toprule
$q$ & $n$ & $\dim V_\lambda$ & $\dim B$ & $B \subseteq E_1^-$ \\
\midrule
$3$ & $10$ & $448$ & $4$    & $\times$ ($T_1$ mixed, non-stable) \\
$4$ & $11$ & $924$ & $4$    & $\times$ ($T_1$ mixed, non-stable) \\
$5$ & $12$ & $1\,728$ & $4$    & \checkmark (predicted; threshold) \\
\bottomrule
\end{tabular}
\end{center}

The $q = 3, 4$ cells confirm the night $T_1 = -1$ paper's prediction
that the $E_1^-$-containment threshold is genuinely sharp at $q_0 = 5$:
$\dim B = 4$ matches the closed-form $f^{(4, 1)} = 4$ already at
$q = 3$, but the $T_1$-action is mixed (non-trivial $+q$-component on
some basis vectors of $B$) until $q \ge 5$. This is the same phenomenon
seen for $(4, 3, 1^q)$ at $q = 4$: dim formula matches but $E_1^-$
fails one step earlier than the structural-threshold prediction.

\paragraph{Phase A decomposition (verified).}
$\dim E_{n-1}^+|_{V_\lambda} = f^{\mu_A} + f^{\mu_B} + f^{\mu_C} +
f^{\nu_1}$ verified at $q = 3$ ($224 = 90 + 35 + 35 + 64$) and
$q = 4$ ($420 = 189 + 70 + 56 + 105$). Direct sum and total
$+q$-eigenspace count of $T_{n-1}|_{V_\lambda}$.

\paragraph{Sub-shapes (sanity).}
\begin{center}
\begin{tabular}{cccc}
\toprule
shape & $\dim V$ & $\dim B$ & $B \subseteq E_1^-$ \\
\midrule
$\mu_A = (4, 2, 1^4)$, $|\mu_A| = 10$ & 1134 & $3$ & \checkmark (Theorem~\ref{thm:muA}) \\
$\mu_B = (5, 1^5)$, $|\mu_B| = 10$    & $126$ & $1$ & \checkmark (Theorem~\ref{thm:muB}) \\
$\mu_C = (4, 1^6)$, $|\mu_C| = 10$    & $84$  & $1$ & \checkmark (Theorem~\ref{thm:muC}) \\
$\nu_1 = (3, 2, 1^5)$, $|\nu_1| = 10$ & $568$ & $2$ & \checkmark (Theorem~\ref{thm:nu1}) \\
\bottomrule
\end{tabular}
\end{center}

\paragraph{Script.}
\texttt{\textasciitilde/projects/scratch/2026-05-13-5-2-1n/verify\_521.py}.

\section{Honest scope and connections}\label{sec:scope}

\paragraph{What is proved structurally at $q \ge 5$.}
$B^{(\lambda)}_{j_0} V_\lambda \subseteq E_1^-|_{V_\lambda}$ and
$\dim \le 4$ for $\lambda = (5, 2, 1^q)$. Rank-zero theorem
$\Pi^{S_n}|_{V_\lambda} = 0$ for $n \ge 12$.

\paragraph{What is conditional.}
None of the structural inputs are open: \textbf{SRD-421},
\textbf{Shift Lemma}, and the May-20 $j$-formula on $(3, 2, 1^q)$ are
all unconditional in their advertised ranges. The four-branch
reduction at $\lambda$ itself relies only on Phase A endpoint
(May-19, unconditional) and the index-gap commutation lemma
(Lemma~\ref{lem:commute}, purely operator-theoretic). So
Theorem~\ref{thm:main} is fully unconditional at $q \ge 5$.

\paragraph{What is deferred to a follow-up.}
The matching structural \emph{lower bound} $\dim B \ge 4$. The
template (disjoint-cell-pair $\Phi$-supports + outer-chain adjacent
injectivity, May-13 \textbf{LB-closed}) applies verbatim here:
$A$- and $B$-cell pairs are $\{c_1, c_3\}$ and $\{c_2, c_3\}$, which
share $c_3$ but differ in the other element ($c_1 \ne c_2$). Hence the
seminormal supports of $\Phi_A(V_{\mu_A})$ and $\Phi_B(V_{\mu_B})$ are
disjoint as sets of SYTs of $\lambda$. Both images lie in
$E_{n-1}^+|_{V_\lambda}$; the outer chain
$\mathcal{O}_\lambda$ is injective on $E_{n-1}^+|_{V_\lambda}$ by
adjacent-injectivity (May-13 \textbf{dim2-331-structural} \S2). The
direct sum has dim $3 + 1 = 4$, preserved by $\mathcal{O}_\lambda$.
The structural inputs (full dim of $\Phi$-images on $\mu_A$, $\mu_B$)
already hold structurally at $q \ge 5$. So the lower bound is
\emph{also} structural at $q \ge 5$, by the same template — promoting
the upper bound to a structural $\dim B = 4$ equality.

\paragraph{What this completes.}
The $|\hat\lambda| = 7$ row of the structural-$j$-formula club:
\begin{center}
\begin{tabular}{ccc}
\toprule
$\hat\lambda$ & shape & structural at \\
\midrule
$(4, 3)$ & $(4, 3, 1^q)$ & $q \ge 5$ (afternoon paper) \\
$(5, 2)$ & $(5, 2, 1^q)$ & $q \ge 5$ (this paper) \\
\bottomrule
\end{tabular}
\end{center}

The two families coincide on threshold ($q_0((5,2)) = q_0((4,3)) = 5$)
but differ on $\dim B$ ($5$ vs.\ $4$, controlled by
$f^{\lambda^{(\ge 2)}}$). Both fit the night $T_1 = -1$ theorem's
threshold formula. The next row $|\hat\lambda| = 8$ contains
$(6, 2), (5, 3), (4, 4), (4, 2, 2), (3, 3, 2)$, etc.\ — the ladder
keeps climbing.

\paragraph{Connection to the closed-form SYT formula.}
With Theorem~\ref{thm:main} and the matching lower bound noted above,
the closed-form SYT formula $\dim B^{(\lambda)} = f^{\lambda^{(\ge 2)}}
= f^{(4, 1)} = 4$ holds structurally at $q \ge 5$ for
$\lambda = (5, 2, 1^q)$, contributing a new structurally verified
$r = 2$ data point to the May-13 evening-1
\texttt{2026-05-13-r-additivity-and-syt-formula.tex} program.

\paragraph{Push status.}
$54 + 1 = 55$ unpushed commits expected on \texttt{clio-vega/proofs}.
PAT still read-only.

\end{document}
